\documentclass[12pt]{article}
\usepackage[T1]{fontenc}
\usepackage[utf8]{inputenc}
\usepackage{lmodern}
\usepackage{textcomp}
\usepackage{enumitem}
\usepackage{geometry}
\geometry{margin=1in}
\begin{document}
\begin{center}
Answer Key - Political Science - Graduate Level Assessment 8 \
\small (Answers are ordered exactly as in the question paper.)
\end{center}

\section*{Multiple Choice}
\begin{enumerate}[label=\arabic*.]
\item \textbf{Answer:} C
\\ \textit{Options:} A) It democratizes countries; B) It has no real effect; C) It encourages stability of the regime; D) It changes the nature of the investor
\\ \textit{Note:} The correct answer is based on the concept of the "resource curse," which suggests that revenues from oil and natural resources can provide regimes in developing states with the financial means to maintain power, suppress dissent, and foster political stability by funding patronage networks and security forces.
\item \textbf{Answer:} C
\\ \textit{Options:} A) The arms trade shifts from being a private to a government controlled enterprise.; B) An increase in the defence trade.; C) A decrease in the arms trade.; D) A growth in the number and variety of weapons traded.
\\ \textit{Note:} The correct answer is based on the fact that the 20 th century witnessed a significant increase in the arms trade due to factors such as global conflict, militarization during the Cold War, and the proliferation of advanced weaponry, making a decrease in the arms trade an inaccurate trend.
\item \textbf{Answer:} D
\\ \textit{Options:} A) James Madison; B) Abraham Lincoln; C) Woodrow Wilson; D) Thomas Jefferson
\\ \textit{Note:} The phrase "Peace, commerce, and honest friendship with all nations, entangling alliances with none" is a direct quote from Thomas Jefferson's inaugural address in 1801, reflecting his foreign policy philosophy of avoiding permanent alliances and promoting neutrality.
\item \textbf{Answer:} D
\\ \textit{Options:} A) Naval victories; B) Diplomacy; C) British preoccupation with Europe; D) All of the above
\\ \textit{Note:} The correct answer "All of the above" signifies that multiple factors, including strategic military decisions, political circumstances, and the socio-economic context of the time, collectively contributed to the United States' ability to avert disaster during the British invasion in 1814, highlighting the complexity of historical events where no single factor can be isolated as the sole reason for the outcome.
\item \textbf{Answer:} C
\\ \textit{Options:} A) It was the first time the US became involved in the Balkans; B) It was the first time NATO used military force; C) It was the first war won by airpower alone; D) It was the first war to employ 'smart weapons'
\\ \textit{Note:} The NATO intervention in Kosovo was unique because it marked the first instance in which a military conflict was successfully concluded solely through aerial bombardment, without the deployment of ground troops, demonstrating the effectiveness of airpower in modern warfare.
\end{enumerate}

\section*{True / False}
\begin{enumerate}[label=\arabic*.]
\setcounter{enumi}{5}
\item \textbf{Answer:} True
\\ \textit{Options:} A) True; B) False
\\ \textit{Note:} The National Security Council is composed of key government officials and advisors who collaborate to formulate and implement the United States' foreign and military policy, making it the primary entity for coordinating these efforts.
\item \textbf{Answer:} True
\\ \textit{Options:} A) True; B) False
\\ \textit{Note:} The term 'indiscriminate attack' implies a lack of precision and intentionality that contradicts the principles of legitimacy and ethical conduct emphasized in the evolving terminology of the Revolution in Military Affairs (RMA), which advocates for precision, discrimination, and minimization of civilian casualties in military operations.
\item \textbf{Answer:} False
\\ \textit{Options:} A) True; B) False
\\ \textit{Note:} The deepening and broadening of security encompasses a range of non-military threats, such as economic instability, environmental challenges, and transnational crime, which extend beyond the traditional focus on military threats posed by sovereign states.
\item \textbf{Answer:} False
\\ \textit{Options:} A) True; B) False
\\ \textit{Note:} The statement is false because while war can be influenced by political objectives, it can also arise from a variety of factors, including economic interests, territorial disputes, cultural conflicts, and humanitarian concerns, which may not necessarily align with partisan political disputes.
\item \textbf{Answer:} True
\\ \textit{Options:} A) True; B) False
\\ \textit{Note:} Bureaucratic politics highlights the intricate decision-making processes within government agencies, which are crucial for effectively identifying, responding to nuclear threats, and managing the complexities of nuclear weapon control.
\end{enumerate}

\section*{Long Form Response}
\begin{enumerate}[label=\arabic*.]
\setcounter{enumi}{10}
\item \textbf{Answer:} The concept of 'empire by invitation' highlights how states may voluntarily seek security from an external power, thus shaping their sovereignty and the dynamics of international relations. This reliance can lead to a reconfiguration of sovereignty, as the host state cedes certain aspects of its autonomy in exchange for protection and stability. Such arrangements can result in a form of dependency that influences domestic policy and limits the host state's ability to act independently on the global stage. Additionally, the invitation of foreign powers often alters regional power balances, prompting responses from neighboring states and potentially leading to new forms of alliance or conflict. Ultimately, while this strategy may enhance immediate security, it raises critical questions about the long-term implications for state sovereignty and the legitimacy of external influence in domestic affairs.
\\ \textit{Note:} The answer correctly identifies 'empire by invitation' as a process where states surrender elements of their sovereignty for external security, illustrating how such dependence alters their autonomy and affects their international relations.
\item \textbf{Answer:} Peace studies in contemporary international relations is characterized by its focus on the underlying causes of conflicts, emphasizing a global perspective that incorporates diverse contexts and cultures. This field is inherently interdisciplinary, drawing from political science, sociology, psychology, and history to analyze the complexities of peace and conflict. A key aspect of peace studies is the pursuit of non-violent transformations, advocating for methods that prioritize dialogue, reconciliation, and social justice over military solutions. Furthermore, it employs an analytical and normative approach that seeks to understand not only the dynamics of conflict but also the ethical implications of peacebuilding efforts. By utilizing multi-level analysis, peace studies effectively links theoretical frameworks with practical applications, ensuring that research contributes meaningfully to real-world peace initiatives.
\\ \textit{Note:} correct because it effectively captures the interdisciplinary nature of peace studies by emphasizing its integration of various academic fields and its focus on understanding the root causes of conflict while advocating for non-violent transformation through dialogue and reconciliation.
\end{enumerate}

\vfill
\noindent\textit{End of Answer Key}
\end{document}